\documentclass[10pt,a4paper]{amsart}
\usepackage{amsmath,amssymb,amsthm,graphicx}
% Retained preamble for the Stacks Project chapter snapshot in this directory.
% The chapter's own preamble.tex is a copy of the shared Stacks Project preamble; the block
% below is sliced MECHANICALLY, line for line, from it. Nothing was added, removed or
% redefined; the wrapper class, the xr-hyper/hyperref wiring and the externaldocument list
% are omitted because the assembler supplies the document class.
%
% lines taken from preamble.tex: 3, 5, 9, 11, 12, 156, 157, 158, ... (51 lines total)
%
% package coverage restored from the chapter's own preamble (2026-10-13): xy, xy-2cell

\usepackage[all]{xy}
\xyoption{2cell}
\UseAllTwocells

\newenvironment{reference}{\comment}{\endcomment}
\newenvironment{slogan}{\comment}{\endcomment}
\newenvironment{history}{\comment}{\endcomment}
\theoremstyle{plain}
\newtheorem{theorem}[subsection]{Theorem}
\newtheorem{proposition}[subsection]{Proposition}
\newtheorem{lemma}[subsection]{Lemma}
\theoremstyle{definition}
\newtheorem{definition}[subsection]{Definition}
\newtheorem{example}[subsection]{Example}
\newtheorem{exercise}[subsection]{Exercise}
\newtheorem{situation}[subsection]{Situation}
\theoremstyle{remark}
\newtheorem{remark}[subsection]{Remark}
\newtheorem{remarks}[subsection]{Remarks}
\numberwithin{equation}{subsection}
\def\lim{\mathop{\mathrm{lim}}\nolimits}
\def\colim{\mathop{\mathrm{colim}}\nolimits}
\def\Spec{\mathop{\mathrm{Spec}}}
\def\Hom{\mathop{\mathrm{Hom}}\nolimits}
\def\Ext{\mathop{\mathrm{Ext}}\nolimits}
\def\SheafHom{\mathop{\mathcal{H}\!\mathit{om}}\nolimits}
\def\SheafExt{\mathop{\mathcal{E}\!\mathit{xt}}\nolimits}
\def\Sch{\mathit{Sch}}
\def\Mor{\mathop{\mathrm{Mor}}\nolimits}
\def\Ob{\mathop{\mathrm{Ob}}\nolimits}
\def\Sh{\mathop{\mathit{Sh}}\nolimits}
\def\NL{\mathop{N\!L}\nolimits}
\def\CH{\mathop{\mathrm{CH}}\nolimits}
\def\proetale{{pro\text{-}\acute{e}tale}}
\def\etale{{\acute{e}tale}}
\def\QCoh{\mathit{QCoh}}
\def\Ker{\mathop{\mathrm{Ker}}}
\def\Im{\mathop{\mathrm{Im}}}
\def\Coker{\mathop{\mathrm{Coker}}}
\def\Coim{\mathop{\mathrm{Coim}}}
\DeclareMathSymbol{\boxtimes}{\mathbin}{AMSa}{"02}
\def\QCohstack{\mathcal{QC}\!\mathit{oh}}
\def\Cohstack{\mathcal{C}\!\mathit{oh}}
\def\Spacesstack{\mathcal{S}\!\mathit{paces}}
\def\Quotfunctor{\mathrm{Quot}}
\def\Hilbfunctor{\mathrm{Hilb}}
\def\Curvesstack{\mathcal{C}\!\mathit{urves}}
\def\Polarizedstack{\mathcal{P}\!\mathit{olarized}}
\def\Complexesstack{\mathcal{C}\!\mathit{omplexes}}
\def\Pic{\mathop{\mathrm{Pic}}\nolimits}
\def\Picardstack{\mathcal{P}\!\mathit{ic}}
\def\Picardfunctor{\mathrm{Pic}}
\def\Deformationcategory{\mathcal{D}\!\mathit{ef}}
\usepackage{cleveref}
\crefname{originallemma}{original lemma}{original lemmas}
\Crefname{originallemma}{Original lemma}{Original lemmas}
\makeatletter
\newlabel{defos-lemma-canonical-class-algebra}{{0DYR}{source}{}{sourcepointer-1}{}}
\newlabel{defos-lemma-canonical-class-algebra@cref}{{[originallemma][0][]0DYR}{[1][0][]source}{}{}{}}
\makeatother
\crefname{originallemma}{original lemma}{original lemmas}
\Crefname{originallemma}{Original lemma}{Original lemmas}
\makeatletter
\newlabel{lemma-lift-equivalence-module-derived}{{0DZ1}{source}{}{sourcepointer-2}{}}
\newlabel{lemma-lift-equivalence-module-derived@cref}{{[originallemma][0][]0DZ1}{[1][0][]source}{}{}{}}
\makeatother
\crefname{originallemma}{original lemma}{original lemmas}
\Crefname{originallemma}{Original lemma}{Original lemmas}
\makeatletter
\newlabel{more-algebra-lemma-base-change-RHom}{{0A6A}{source}{}{sourcepointer-3}{}}
\newlabel{more-algebra-lemma-base-change-RHom@cref}{{[originallemma][0][]0A6A}{[1][0][]source}{}{}{}}
\makeatother
\crefname{originallemma}{original lemma}{original lemmas}
\Crefname{originallemma}{Original lemma}{Original lemmas}
\makeatletter
\newlabel{more-algebra-lemma-neighbourhood-equivalence}{{05EC}{source}{}{sourcepointer-4}{}}
\newlabel{more-algebra-lemma-neighbourhood-equivalence@cref}{{[originallemma][0][]05EC}{[1][0][]source}{}{}{}}
\makeatother
\title{Stacks, Examples of Deformations: lemma-lift-equivalence-module}
\author{The Stacks Project Authors}
\date{stacks-project:examples-defos.tex:local-snapshot-20261012; GNU FDL 1.2 or later}
\begin{document}
\maketitle
\noindent\textit{Editorial note.} Statement span is lines 2212-2235 of sources/stacks-examples-defos/pinned-original-source.tex; proof span is lines 2237-2290, the complete author proof (52 body lines).\par\medskip
\noindent\textit{Editorial note.} Navigation stubs: 4 cited-but-not-included labels are supplied as non-load-bearing pointers with their pinned tags where the snapshot tag index has them.\par\medskip
\noindent\textit{Editorial note.} extra\_packages is empty on purpose: the retained preamble already loads xy with options and a second declaration clashes in the assembler wrapper.\par\medskip
\noindent\textit{Editorial note.} Text is retained verbatim, including the author's own typos and placeholders.\par\medskip
\noindent\textit{Editorial note.} Queue index 7 of queue\_examples-defos.json (pb=2202, stmt 2212-2235, proof 2237-2290).\par\medskip
\noindent\textit{Allowed standard prerequisites (editorial).}\begin{itemize}
\item Explicit prerequisite boundary. The cited inputs are prerequisites; the argument in the retained span is not replaced by any cited result.
\end{itemize}\medskip
\ifdefined\hypertarget\hypertarget{sourcepointer-1}{}\fi
\noindent\textit{Source reference.} Original lemma 0DYR: sources/stacks-examples-defos/pinned-original-source.tex, source lines 2290--2290. Cited by the retained statement or proof but not part of the collected spans; supplied only as a navigation pointer. The argument does not depend on its content.\par\medskip
\ifdefined\hypertarget\hypertarget{sourcepointer-2}{}\fi
\noindent\textit{Source reference.} Original lemma 0DZ1: sources/stacks-examples-defos/pinned-original-source.tex, source lines 2170--2181. Cited by the retained statement or proof but not part of the collected spans; supplied only as a navigation pointer. The argument does not depend on its content.\par\medskip
\ifdefined\hypertarget\hypertarget{sourcepointer-3}{}\fi
\noindent\textit{Source reference.} Original lemma 0A6A: sources/stacks-examples-defos/pinned-original-source.tex, source lines 2290--2290. Cited by the retained statement or proof but not part of the collected spans; supplied only as a navigation pointer. The argument does not depend on its content.\par\medskip
\ifdefined\hypertarget\hypertarget{sourcepointer-4}{}\fi
\noindent\textit{Source reference.} Original lemma 05EC: sources/stacks-examples-defos/pinned-original-source.tex, source lines 2290--2290. Cited by the retained statement or proof but not part of the collected spans; supplied only as a navigation pointer. The argument does not depend on its content.\par\medskip
\begin{lemma}
\label{lemma-lift-equivalence-module}
Consider a commutative diagram of Noetherian rings
$$
\xymatrix{
A' \ar[d] \ar[r] &
P' \ar[d] \ar[r] &
Q' \ar[d] \\
A \ar[r] &
P \ar[r] &
Q
}
$$
with cartesian squares, with flat horizontal arrows, and with
surjective vertical arrows whose kernels are nilpotent.
Let $J' \subset P'$ be an ideal such that $P'/J' = Q'/J'Q'$.
Let $M$ be an $A$-flat $P$-module.
Assume for all $g \in J'$ there exists an $A'$-flat $(P')_g$-module
lifting $M_g$. Then the following are equivalent
\begin{enumerate}
\item $M$ has an $A'$-flat lift to a $P'$-module, and
\item $M \otimes_P Q$ has an $A'$-flat lift to a $Q'$-module.
\end{enumerate}
\end{lemma}

\begin{proof}
Let $I = \Ker(A' \to A)$. By induction on the integer $n > 1$
such that $I^n = 0$ we reduce to the case where $I$ is an ideal
of square zero; details omitted.
We translate the condition of liftability of
$M$ into the problem of finding an object of $D^-(P')$ as in
Lemma \ref{lemma-lift-equivalence-module-derived}.
The obstruction to doing this is the element
$$
\omega(M) \in \text{Ext}^2_P(M, M \otimes_P^\mathbf{L} IP) =
\text{Ext}^2_P(M, M \otimes_P IP)
$$
constructed in
Deformation Theory, Lemma \ref{defos-lemma-canonical-class-algebra}.
The equality in the displayed formula holds as
$M \otimes_P^\mathbf{L} IP = M \otimes_P IP$
since $M$ and $P$ are $A$-flat\footnote{Choose a resolution
$F_\bullet \to I$ by free $A$-modules. Since $A \to P$ is flat,
$P \otimes_A F_\bullet$ is a free resolution of $IP$.
Hence $M \otimes_P^\mathbf{L} IP$ is represented by
$M \otimes_P P \otimes_A F_\bullet = M \otimes_A F_\bullet$.
This only has cohomology in degree $0$ as $M$ is $A$-flat.}.
The obstruction for lifting $M \otimes_P Q$ is similarly
the element
$$
\omega(M \otimes_P Q) \in
\text{Ext}^2_Q(M \otimes_P Q, (M \otimes_P Q) \otimes_Q IQ)
$$
which is the image of $\omega(M)$ by the functoriality
of the construction $\omega(-)$ of
Deformation Theory, Lemma \ref{defos-lemma-canonical-class-algebra}.
By More on Algebra, Lemma \ref{more-algebra-lemma-base-change-RHom}
we have
$$
\text{Ext}^2_Q(M \otimes_P Q, (M \otimes_P Q) \otimes_Q IQ) =
\text{Ext}^2_P(M, M \otimes_P IP) \otimes_P Q
$$
here we use that $P$ is Noetherian and $M$ finite.
Our assumption on $P' \to Q'$ guarantees that for an $P$-module $E$
the map $E \to E \otimes_P Q$ is bijective on $J'$-power torsion, see
More on Algebra, Lemma
\ref{more-algebra-lemma-neighbourhood-equivalence}.
Thus we conclude that it suffices to show $\omega(M)$
is $J'$-power torsion. In other words, it suffices to show that
$\omega(M)$ dies in
$$
\text{Ext}^2_P(M, M \otimes_P IP)_g =
\text{Ext}^2_{P_g}(M_g, M_g \otimes_{P_g} IP_g)
$$
for all $g \in J'$. Howeover, by the compatibility of formation of $\omega(M)$
with base change again, we conclude that this is true as $M_g$
is assumed to have a lift (of course you have to use the whole
string of equivalences again).
\end{proof}

\end{document}
